\section{Scope of the Criterion}\label{app:boundaries}
\subsection{Other Dissimilarities}
The shrinking-pair statement gives a criterion-transfer corollary over the full continuous competitor class.

\begin{corollary}
Suppose $D:[0,\infty)\to[0,\infty)$ is continuous, $D(0)=0$, positive away from zero, and
$$
 D(d)\ge c_Dd^2\quad(0\le d\le2K\sqrt P),\qquad
 \lim_{d\downarrow0}\frac{D(d)}{d^2}=c_D>0.
$$
For the targets of Proposition~\ref{prop:score}, every exact continuous encoder satisfies $\Lambda_D=c_D\Lambda_{\sq}$. Therefore every optimal encoder for the targets of Proposition~\ref{prop:score} remains optimal, and its optimum is multiplied by $c_D$.
\end{corollary}
\begin{proof}
The lower envelope gives $\Lambda_D\ge c_D\Lambda_{\sq}$. Proposition~\ref{prop:score} supplies pairs attaining the opposite inequality in the limit. Nonexact encoders have zero score because $D(0)=0$.
\end{proof}

Examples are $D(d)=d^2+ad^4$ with $a\ge0$ and $D(d)=e^{ad^2}-1$ with $a>0$. The original uncentered $e^{ad^2}$ criterion has $D(0)=1$ and is not covered. No closed-form optimizer for that criterion is claimed here. Gaussian MSE and optimality at a fixed disturbance radius also remain separate problems.


\subsection{Exactness and Conditioning}
The family
$$\enc_\beta(s)=\sqrt P[s+i\beta(1-s^2)],\qquad0<\beta\le1,$$
computes $mq$ exactly, but its squared score is
$$\Lambda_\sq(\enc_\beta;mq)=K^2P\frac{\beta^2}{1+\beta^2}.$$
Its inverse conditioning therefore diverges as $\beta\downarrow0$. At zero and $K=4$, inputs $(1,1,0,0)$ and $(1/2,1/2,1/2,1/2)$ give the same received sum while their targets are $1/4$ and $1/8$.

For a nonconstant bounded target, the uncentered exponential criterion assigns every encoder a positive score of at least $1/(\max f-\min f)^2$, even if exactness fails. Its nonzero value at zero distance therefore makes exactness a separate condition. Noise-dependent exponential optimization requires a separate analysis.

\subsection{Finite Alphabets}
For $K=4$ and $s_k\in\{-1,1\}$, the possible means are $m\in\{-1,-1/2,0,1/2,1\}$ and variance is $1-m^2$. Antipodal signaling gives received sum $4\sqrt P m$, hence
$$\frac{|\Delta r|^2}{|\Delta V|^2}=\frac{16P}{(m+\widetilde m)^2}$$
whenever the target gap is nonzero. The largest admissible $|m+\widetilde m|$ is $3/2$, giving score $64P/9$. Every binary map has symbol separation at most $2\sqrt P$, so antipodal signaling maximizes every received gap and is globally optimal. Restricting the continuous variance-optimal parabola instead gives
$$\left(\frac4{\sqrt3}-\frac43\right)\frac{64P}{9}<\frac{64P}{9}.$$
There is no contradiction: removing source values removes the continuum's sharp shrinking pairs.

